\section{Where the wall is}\label{sec:wall}

This section is a diagnosis. It contains no new mathematics. Its claims are consequences of the results cited in
\cref{sec:series}, and where it gives a reading instead of a consequence it says so.

\subsection{What is left is one function}\label{sec:one-function}

For oblivious machines the series has reduced~(T) to a question about codes (\cref{sec:depth}). Fix the number of
tapes, the alphabet and the time constant, and let $\eps>0$. Depth reduction gives, for every linear-time schedule, a
genie that reveals at most $\eps L$ bits and leaves cones of depth at most $D_\eps$, a bound that depends on $\eps$ and
on the fixed parameters but not on~$d$. Write $r(\eps)$ for the best rate, in the limit of large~$d$, of a binary
storage code on the cone relations that such genies leave on such schedules. Then:
\begin{itemize}[leftmargin=1.5em]
\item an oblivious linear-time transposer needs $r(\eps)\ge1-\eps$ for every $\eps>0$;
\item a proof that $r(\eps)<1-\eps$ for one $\eps$, for each choice of the fixed parameters, would exclude oblivious
linear-time transposers~\cite[Corollary~9.3]{BoundedStack2026};
\item the pairing gives $r(\eps)\ge\tfrac12$, and whether $r(\eps)>\tfrac12$ for any~$\eps$ is open in both directions.
\end{itemize}
The depth bound matters: without it the cone of an output is everything read before it, and the question is empty.
The only general upper bound is the elementary one for relations of degree at most~$\Delta$, a rate of
$1-1/(\Delta+1)$~\cite[Proposition~4.14]{FourthPaper2026}, which is useless at the depths in question. For machines with linear data paths the same reduction
ends in the relative rigidity of the transposition matrix on ball patterns (\cref{sec:depth}). For adaptive machines
the question arises in two forms: the Tape Coding Conjecture of the first paper, which needs no simulator, and
the storage-code criteria on the traces of a bit simulator (\cref{sec:adaptive}).

For particular classes of schedules the exclusion is proved. The sixteenth paper proves it for schedules whose
large output sets can be cut cheaply from the macroscopic crossings, and for schedules with boundedly many macroscopic
phases~\cite[Theorems~5.4, 9.2 and~9.3]{SixteenthPaper2026}; the source covers of the ninth paper exclude restricted
geometries for all codes. What is missing is a statement for all schedules.

\subsection{Three kinds of argument and their limits}\label{sec:exhaustions}

\paragraph{Motion without data.} The cut inequality behind Theorem~M*, its entry-cell refinement and every conclusion
of Theorem~M* are satisfied by the padded machine of the second paper, which transposes nothing (\cref{sec:motion}).
So conditions on head motion that do not refer to the outputs are necessary and cannot be decisive. Conditions that
do refer to them go further: the rows-to-columns inequality of the sixteenth paper excludes two of the five candidate
schedules. But one cut serves every pair of row sets, so this inequality gives a logarithm only when the cuts for
different address bits are disjoint, which is a hypothesis and not a theorem; and what excludes all five candidates is
the regular-recursion obstruction, which applies only to schedules all of whose stages have regular limiting motion.
Whether conditions of this kind decide the oblivious case is open.

\paragraph{Codes without machines.} On abstract separated relations of bounded degree, storage codes exceed the pairing
rate, and their rates tend to one polynomially fast as the degree cap grows (\cref{sec:false}). A transposer is only
shown to need rates that tend to one very slowly, and the cap grows as $\eps$ falls because a genie of cost $\eps L$
has depth at least of order $1/\eps$~\cite[Proposition~4.13]{FourthPaper2026}. So no rate bound for separated relations of bounded
degree in general can separate a transposer from an abstract code. The two properties of cone relations that the
series' coding bounds use, bounded degree after trimming and paired-distance separation, are shared by the
counterexamples.

\paragraph{Counting passes.} The pass bound of the first paper is the one logarithm proved for a resource of every
machine. It saturates one scale at a time: a linear-time run has room for about $2^j$ legs at scale $2^{-j}L$, while
counting at a single scale asks for only $O(j)$ of them, and the threshold of the coarse-pass bound of the sixteenth
paper is tight~\cite[Theorem~10.3 and Corollary~10.4]{SixteenthPaper2026}. That paper states the problem of an
inequality in the numbers of legs at all scales at once~\cite[Problem~11.1]{SixteenthPaper2026}; nothing in the series
proves one.

\subsection{Why every general tool saturates}\label{sec:saturation}

The $\log d$ of the folklore algorithm is a sum over $\log_2d$ address digits. Exchanging one digit of the row address
with one digit of the column address costs $O(L)$, and the transposition is the product of these $\log_2d$ commuting
exchanges. Linear time pays for any one of them in full, and for any bounded number.

Every tool of the series \emph{that applies to all machines} is single-scale or scale-free. Cut arguments, macroscopic
limits and coarse passes look at one scale at a time, and a linear-time machine has enough steps to satisfy any
single-scale requirement. Storage-code bounds are scale-free: they see a relation, not the scale of the addresses that
produced it. The one attempt to make a geometric quantity grow with the number of scales failed: walks nested at $m$
scales are cut to a depth independent of~$m$ (\cref{sec:depth}).

The series does charge scales separately in restricted settings, and they show what the general case lacks.
\begin{itemize}[leftmargin=1.5em]
\item For transport at bounded occupancy, each moved label is charged its $\log d$ address bits, in
steps~\cite[Lemma~3.4]{DirectedRecovery2026}. What the bound does not cover is coding.
\item For traces whose cuts for different address bits are disjoint, the rows-to-columns inequality gives
$\Omega(L\log d)$ against every code~\cite[Corollary~7.3]{SixteenthPaper2026}. What is not available is the
disjointness.
\item In the abstract model of XOR-frame recursion every level is charged~\cite{FirstPaper2026}, and compilers with
boundedly many stages meet a multiscale obstruction~\cite{ThirdPaper2026}. Both are restricted models.
\end{itemize}
So the logarithm is proved wherever either the coding or the geometry is tied down, and nowhere else.
\Cref{cor:digit,rem:compose-TP} below make one half of this precise in the linear one-hop model: every address digit
must be treated asymmetrically, but one cheap tape operation does that for all digits at once.

\subsection{What the missing statement must be}\label{sec:missing}

The limits above describe the statement that would decide the oblivious case. It must
\begin{itemize}[leftmargin=1.5em]
\item tie what is known to where it can be, for all schedules: the class results of \cref{sec:one-function} do this
for restricted geometries;
\item fail for abstract separated relations, where codes of rate tending to one exist;
\item fail for motion that ignores the data, which padding supplies;
\item constrain all scales at once;
\item and use the shape of tapes, since a bound on coding gains for bounded-degree graphs in general would be a circuit
lower bound (\cref{sec:first}).
\end{itemize}
Since a linear-time multiplier gives a linear-time transposer~\cite{HvdH2025}, a proof of such a statement for all
machines would also be the first superlinear lower bound for off-line multiplication on multitape machines.

\subsection{What a construction cannot avoid}\label{sec:needs}

The same results say what a linear-time machine would have to use.
\begin{itemize}[leftmargin=1.5em]
\item \emph{No bounded complex coefficients.} Linear recovery with bounded complex coefficients and bounded width
stays at the pairing rate on cone relations (\cref{sec:proved}). A linear code would need positive characteristic or
coefficients that cannot be rescaled to be bounded; nonlinear codes are not covered.
\item \emph{Curved motion.} In the limit, head motion must fail to be locally affine on a positive fraction of the
time (\cref{sec:motion}), and a recursive schedule must have a stage whose limiting motion is not regular.
\item \emph{Expanding relations.} The trimmed incidence graphs of the cone relations must not be uniformly hyperfinite
(\cref{sec:proved}). This is necessary and cheap: a fixed linear-time machine has such cones and still obeys the
all-code half bound~\cite[Corollary~7.9]{DirectedRecovery2026}.
\item \emph{No gadgets.} After trimming, almost no output is within bounded distance of its own input bit, so no gain
above the pairing rate can come from gadgets of bounded size (\cref{sec:proved}). The pairing itself is not excluded;
it stops at~$\tfrac12$.
\end{itemize}
Above the pairing rate, at a degree cap that is fixed before the length and the paired distances grow, the only codes
known come from one source: sections of line bundles on the supersingular loci of Hilbert modular
surfaces~\cite{StorageCodes2026}, by a method that originates in~\cite{Luck2026}, together with their
amplifications~\cite{HighRate2026}. At degrees that grow with the length other examples exist.

\subsection{The rule and the relation}\label{sec:rule}

One feature of machines is absent from the coding statements above. A storage code on a relation lets each receiver
choose its own decoding function. A Turing machine has one finite transition table, and on a stretch of time during
which its heads move in step, the function that produces the output at one position is the function that produces the
output at the next position, translated. The cone relation records which inputs each output can depend on. It does not
record that neighbouring outputs depend on them \emph{in the same way}.

Sharing decoders is not by itself an obstruction to rate above one half: by their construction, the linear codes of
the thirteenth paper are equivariant under large groups with boundedly many orbits of coordinates. But transposition has a feature that an
abstract relation lacks: its demand is not invariant under translating one index. Shifting the columns of a matrix
shifts the rows of its transpose. The two theorems of this paper are the two cleanest ways in which structure of the
decoder on one index side is incompatible with the swap: \emph{order} (\cref{sec:order}) and \emph{symmetry}
(\cref{sec:symmetry}). They are proved in the linear one-hop model of \cref{sec:onehop}. They are not theorems about
machines, and \cref{sec:construction} says what would be needed to make them so.
